\scanpage{33}

Catégorie de Picard-type non stricte.

$A$ anneau commutatif (avec 1)

\deleted{\illegible{2}} $2.1_A\ne 0$

\[
C\quad\left|\begin{aligned}
\operatorname{Ob} C={}&\text{Modules gradués \inserted{mod 2} sur \(A\) munis d'une \uline{base} formée}\\
&\text{d'un élément, mais donnée seulement}\\
&\text{au signe près \deleted{\unsure{à base}}\unsure{(équivt.|(partie|[partie|(soit} : deux éléments}\\
&\text{\(e_L\), stable par \(x\longmapsto-x\), formée d'\unsure{éléments} bases)}\\
\operatorname{Hom}(L,M)={}&\text{Ensemble des isomorphismes \inserted{gradués} \(\varphi:L\longrightarrow M\) tels}\\
&\text{que \(\varphi(e_L)=e_M\)}\\
\otimes={}&\text{produit tensoriel gradué, avec}\\
& e_{(L\otimes M)}=\{x\otimes y\mid x\in e_L,\;y\in e_M\}
\end{aligned}\right.
\]

donnée d'associativité : comme d'habitude \unsure{des} (\unsure{on peut typer la structure des bases|on peut typer la structure des bibases|compatible d'après la structure des bases})

donnée de commutativité : \uline{signe} de Koszul

\uline{description} \uline{alternative}
\[
\deleted{\unsure{\mathcal C|C}}C=
 \underbrace{\left(\underbrace{\mathbb Z/2\mathbb Z}_{\substack{\text{Catégorie}\\\text{discrète}}}\right)
 \times\underbrace{C_0}_{\substack{\longrightarrow\,\text{catégorie des}\\\text{torseurs sous}\\\mathbb Z/2\mathbb Z}}}
 _{\substack{\deleted{\unsure{\otimes}}\ \text{torseurs gradués mod}\\2\ \text{sous }\mathbb Z/2\mathbb Z}}
\]

\uline{contraintes} ass : comme \unsure{habitude}

\hspace*{3em}comm : \deleted{\unsure{par}} règle \deleted{\unsure{K}}Koszul

\deleted{\uline{Torseurs gradués}}

\deleted{\(F\) faisceau abélien sur topos \(S\) (écrit additivement)}

\deleted{\(\Gamma\) groupe discret}

\[
\deleted{\Gamma\times\Gamma\xrightarrow{f}H^0(S,F)\qquad\text{application}}
\]

\[
\deleted{C=\underbrace{H^{0}(S,\Gamma_S)}_{\substack{\text{catégorie}\\\text{discrète}\\\text{d'ens. ss-jacent}\\H^0(S,\Gamma_S)}}
\times\underbrace{\underline{\operatorname{Tors}}(F)}_{\substack{\text{catégorie des}\\\text{torseurs sous \(F\)}}}}
\]

\diagram{\begin{tikzcd}[ampersand replacement=\&]
 \deleted{C\times C}\arrow[r,"\deleted{\otimes}"]\&\deleted{C}
\end{tikzcd}}
\[
\deleted{(\alpha,P)\otimes(\beta,Q)=(\alpha\beta,P\otimes Q)}
\]
\deleted{associativité}
\diagram{\resizebox{\linewidth}{!}{\begin{tikzcd}[ampersand replacement=\&,row sep=1.3ex,column sep=1em]
 \deleted{\big((\alpha,P)\otimes(\beta,Q)\big)\otimes(\gamma,R)}\arrow[d,equal]\&
 \deleted{\simeq}\&\deleted{(\alpha,P)\otimes\unsure{\big((\beta,Q)\otimes(\gamma,R)\big)|\big(\beta,6,\otimes\gamma R\big)}}\arrow[d,equal]\\
 \deleted{\big(\alpha\beta\gamma,(P\otimes Q)\otimes R\big)}\& \&
 \deleted{\big(\alpha\beta\gamma,P\otimes(Q\otimes R)\big)}
\end{tikzcd}}}

\deleted{provient de associativité de \(\otimes\) dans \uline{Tors} \((F)\)}
% archive sheet number: 32
